\begin{enumerate}
	\item Se cumplió el primer objetivo específico. Se rediseñaron en Autodesk Inventor los carros, el conjunto central y el mecanismo vertical para incorporar ejes de acero de 12~mm y rodamientos lineales, y se fabricó el mecanismo con 48 piezas en PLA. Después del montaje no se reprodujeron los atascos del sistema original y el avance desigual del pórtico se volvió poco perceptible a simple vista; la lubricación del sistema original no había corregido estas fallas. La mejora fue cualitativa, ya que no se midieron fuerzas de desplazamiento ni deflexiones.

	\item Se cumplió el segundo objetivo específico. El modo manual permitió mover X, Y y Z de forma independiente, orientar el efector final y accionar la pinza desde un mando inalámbrico, con bloqueo del sentido de cada final de carrera activo, neutralización de las órdenes al perder el mando y regreso al origen desde el propio mando.

	\item Se cumplió parcialmente el tercer objetivo específico. Se integraron la cámara HUSKYLENS 2 con un modelo YOLO11n de segmentación, una homografía calibrada con cuatro AprilTags, seis finales de carrera y un encoder en la banda, y el sistema recogió piezas en movimiento y las trasladó a un punto de entrega. El encoder cerró el lazo a nivel de supervisión, al actualizar la posición de la pieza con el avance medido de la banda; la posición de los ejes del brazo, en cambio, se controló en lazo abierto mediante el conteo de pasos.

	\item Se cumplió parcialmente el cuarto objetivo específico. En 120 ciclos automáticos, el 61,7~\% de las piezas se recogió y se entregó correctamente (IC 95~\%: 52,7--69,9~\%). La tasa fue mayor con la banda al 25~\% (73,3~\%) que al 35~\% (50,0~\%), con la pieza 7 (75,0~\%) que con la pieza 6 (48,3~\%) y con las piezas verticales (90,6~\%) que con las horizontales (45,2~\%) y las diagonales (54,4~\%). No se midieron la exactitud ni la repetibilidad de pose del efector final según ISO~9283.

	\item La referenciación por finales de carrera mostró desviaciones estándar del recorrido medido de 1,29~mm en X y 0,46~mm en Y, y el modelo de posición de la pieza alcanzó un error al contacto de $-1\pm24$~mm en 22 agarres manuales. Estos indicadores acotaron las fuentes de error del sistema, pero no sustituyen la medición de exactitud y repetibilidad del efector final.

	\item La migración del enlace entre la Portenta H7 y la ESP32 de I2C a UART sobre RS485 redujo los errores de comunicación: las tasas registradas bajaron de 10--16~\% con I2C a menos de 2~\% con RS485, aunque en condiciones de prueba distintas y sin una serie final de funcionamiento continuo. La capa de transporte con COBS, CRC16, sesión y número de solicitud permitió detectar tramas dañadas, descartar respuestas ajenas y recuperar el enlace sin perder las calibraciones.

	\item La verificación del recorrido de las piezas con el encoder permitió detectar y corregir errores de escala y de signo en la calibración de la cámara que se compensaban entre sí. Esta experiencia mostró que una homografía con solución numérica válida no garantiza una transformación exacta, y que una referencia independiente es necesaria para validarla.

	\item La principal limitación del ciclo automático fue la información que entrega la cámara: al recibir solo la caja delimitadora, el sistema no pudo determinar el sentido de inclinación de las piezas diagonales, y la orientación dependió de la clase detectada. Junto con el ajuste empírico del instante de captura, esta limitación es una explicación posible de las menores tasas en piezas horizontales y diagonales.
\end{enumerate}
